\chapter{Conclusion and Future Work}

\section{Conclusion}

This project presented NLGCN, a novel learning-based framework for
influential node detection in weighted graphs \cite{guo2026}. By
combining local and global structural features, Node Local Influence
(NLI) and Node Global Influence (NGI), at multiple aggregation
scales and encoding them as a multi-channel tensor representation,
the model captures rich structural information that traditional
centrality measures such as degree centrality,
K-Shell decomposition, and betweenness centrality
 fail to exploit.

A channel attention mechanism enables the model
to dynamically weight the relative importance of each structural
channel, adapting its representation to different network types
without manual feature selection. Ground truth labels derived from
the SIR epidemic diffusion model provide a
realistic and simulation-validated measure of node influence. The
trained NLGCN model achieves Kendall Tau correlations of 0.8944 to
0.9627 across four real-world benchmark networks, significantly
outperforming all baseline methods, and
achieves perfect top-10 influential node identification on every
dataset evaluated.

The implementation demonstrates that combining domain-specific
structural feature engineering with convolutional deep learning
\cite{zhang2019gcn} yields a fast, scalable, and accurate
approach to influence ranking, replacing expensive per-node SIR
simulations at inference time with a single forward pass through
the trained network \cite{guo2026, patel2025}.

\section{Future Work}

Several directions are identified for extending and strengthening
the NLGCN framework. The most significant planned enhancement
concerns the incorporation of edge weights, which is described
first and in greatest detail.

\subsection{Extension to Weighted Graphs}

The current implementation treats all graphs as unweighted and
undirected, representing edges as binary entries in the adjacency
matrix \cite{guo2026}. This is a significant limitation, as most
real-world networks are inherently weighted. In transportation
networks, edge weights may represent distances or travel times;
in social networks, they may encode interaction frequency or
relationship strength; in biological networks, they may reflect
the strength of synaptic connections or protein interaction
confidence scores \cite{hafiene2020, thotakura2024}.

Ignoring edge weights discards potentially critical information
about node influence. A node connected to a few strong links
may be more influential than one with many weak connections a distinction that the current binary adjacency matrix
cannot capture \cite{mukhtar, chen2012}.

Extending NLGCN to weighted graphs requires modifications at
multiple stages of the pipeline \cite{guo2026, patel2025}:

\begin{itemize}
    \item \textbf{Weighted adjacency matrix:} Replace the binary
    adjacency matrix $A_{ij} \in \{0, 1\}$ with a weighted matrix
    $W_{ij} \in \mathbb{R}_{\geq 0}$, where $W_{ij}$ encodes the
    edge weight between nodes $i$ and $j$ \cite{zhang2019gcn}.

    \item \textbf{Weighted NLI:} The Node Local Influence formula
    should incorporate edge weights into the neighbour contribution
    term, replacing the binary existence of an edge with a
    weighted influence score. Specifically, the exponential
    aggregation over neighbours could be modified to:
    \begin{equation}
        \text{NLI}_i^{w} = \frac{d_w(v_i) \cdot
        \log\!\left(\sum_{j \in N(i)} w_{ij} \cdot
        \exp\!\left(d(v_j)\right)\right)}{n}
    \end{equation}
    where $d_w(v_i) = \sum_{j \in N(i)} w_{ij}$ is the weighted
    degree (node strength) of node $i$ \cite{guo2026, mukhtar}.

    \item \textbf{Weighted NGI:} The Node Global Influence formula
    should use weighted shortest path distances rather than
    hop-count distances, computed via Dijkstra's algorithm on
    the weighted graph \cite{farooq2018, chen2012}:
    \begin{equation}
        \text{NGI}_i^{w} = \sum_{i \neq j}
        \frac{\sqrt{d_w(v_j) + \alpha}}{d_{ij}^{w}}
    \end{equation}
    where $d_{ij}^{w}$ is the weighted shortest path distance
    between nodes $i$ and $j$ \cite{guo2026}.

    \item \textbf{Weighted neighbourhood matrix:} The neighbourhood
    matrix used for channel construction should embed edge weights
    in the off-diagonal entries, replacing the binary $\{0, 1\}$
    encoding with the actual weight $w_{ij}$ \cite{tu2025, patel2025}.
    This directly encodes link strength into the spatial
    representation that the convolutional layers process.

    \item \textbf{Weighted SIR labels:} The label generation
    procedure should also reflect edge weights, with infection
    probability $\beta_{ij} = \beta \cdot w_{ij}$ varying per
    edge rather than being uniform across the network.
\end{itemize}

These modifications would allow NLGCN to fully exploit the
information encoded in edge weights, likely yielding further
improvements in ranking accuracy on real-world weighted networks
\cite{guo2026, zhang2019gcn}.

\subsection{Extension to Directed and Dynamic Graphs}

The current framework assumes undirected, static graphs
\cite{guo2026}. Real-world networks such as Twitter follower
graphs, citation networks, and web graphs are directed, and many
networks including social media platforms and transportation
systems evolve over time \cite{hafiene2020}.



\subsection{Exploring Alternative Diffusion Models}

The current implementation uses the SIR epidemic model exclusively for ground truth label generation.
While SIR is widely used in the influential node detection
literature \cite{guo2026, chen2012}, other diffusion
models such as the Independent Cascade (IC) model and the
Linear Threshold (LT) model \cite{hafiene2020} capture
different spreading dynamics and may be more appropriate for
certain application domains such as viral marketing or opinion
formation.

\subsection{Scalability to Large Graphs}

As discussed in Section~4.2, the current implementation becomes
computationally expensive for very large graphs due to the
all-pairs shortest path computation required for NGI
\cite{guo2026, farooq2018}. Future work should explore
approximate distance methods such as landmark-based shortest
path estimation, graph sampling strategies similar to
GraphSAGE, or hierarchical neighbourhood
construction to reduce the $O(nm)$ complexity of distance
computation while preserving ranking accuracy.

